% WHAT \noexpand HOLDS BACK, AND FOR HOW LONG.
%
% \noexpand holds the name behind it back for ONE READING. The reading that
% spends the holding gets a name that behaves as \relax -- so the name is
% OBEYED, does nothing, and is TRACED:
%
%   \noexpand\undefinedtwo      {\relax}      and no fault at all
%
% HSTeX dropped the name without a word, so \tracingcommands lost a line
% wherever a document writes \noexpand outside a text being gathered.
%
% \expandafter READS the name it holds back, and that reading is the one
% that spends the holding: what goes back is the bare name, expanded when
% it is read again. So
%
%   \expandafter\expandafter\noexpand\undefined
%
% -- where the inner \expandafter fetches the name \noexpand had just held
% -- draws
%
%   ! Undefined control sequence.
%   <recently read> \undefined
%
% where a plain `\noexpand\undefined' draws nothing. trip line 437 writes
% exactly that pair.
%
% A name that was never expandable is not held back at all: \noexpand\relax
% is \relax, and a character token passes through untouched.
\catcode`\{=1 \catcode`\}=2
\tracingonline=1 \tracingcommands=2
\message{[A the expandafter-expandafter-noexpand idiom over an undefined name]}
\expandafter\expandafter\noexpand\undefined\noexpand\expandafter\relax
\message{[B noexpand alone over an undefined name]}
\noexpand\undefinedtwo
\message{[C noexpand over a defined macro]}
\def\m{X}\setbox0=\hbox{\noexpand\m}\showthe\wd0
\message{[D noexpand over an unexpandable primitive]}
\noexpand\relax
\message{[E noexpand over a character]}
\setbox0=\hbox{\noexpand A}\showthe\wd0
\message{[done]}
\end
